% Source: https://artofproblemsolving.com/community/c3h1215412_algebra_help
% Stable identity: aops-topic-1215412-q01
Both inequalities are true, but the second one is not sharp: $t=1$ would need $(a-b)^2+(a+b)^2=0$, that is $a=b=0$, which violates $a^2+ab+b^2=1$. Showing that $t$ lies in an interval does not show that every value of the interval occurs.

Let $s=a^2+b^2$ and $p=ab$. Then $s+p=1$ and $t=p-s=2p-1$. For real $a,b$ we have $s\ge 2|p|$ (since $(a\pm b)^2\ge0$), and conversely any $s\ge2|p|$ comes from real $a,b$. With $s=1-p$:
\begin{itemize}
\item if $p\ge0$: $1-p\ge2p$ gives $p\le\dfrac{1}{3}$;
\item if $p<0$: $1-p\ge-2p$ gives $p\ge-1$.
\end{itemize}
So $-1\le p\le\dfrac{1}{3}$ and $t=2p-1$ ranges over $\boxed{-3\le t\le-\dfrac{1}{3}}$. Both ends occur: $a=1$, $b=-1$ gives $t=-3$, and $a=b=\dfrac{1}{\sqrt3}$ gives $t=-\dfrac{1}{3}$; every value in between occurs because $p$ varies continuously.
